\documentclass[multi=tikzpicture,border=9pt]{standalone}
\usepackage[UTF8]{ctex}
\usepackage{amsmath,amssymb}
\usepackage{tikz}
\usetikzlibrary{arrows.meta,calc,patterns,decorations.pathmorphing}
\definecolor{main}{HTML}{17365D}
\definecolor{accent}{HTML}{AC3931}
\tikzset{every picture/.style={font=\small,line width=.75pt},
>={Stealth[length=2mm]},arr/.style={->},
box/.style={draw=main,rounded corners=2pt,fill=main!4,text width=3.6cm,minimum height=1cm,align=center},
note/.style={align=center,text width=9cm}}
\begin{document}

\begin{tikzpicture}
\draw (0,3)--(.6,3);\draw[decorate,decoration={coil,aspect=.5,segment length=5pt,amplitude=4pt}] (.6,3)--(2.6,3);
\draw (2.6,3)--(3,3)--(3.4,3);\draw[decorate,decoration={coil,aspect=.5,segment length=5pt,amplitude=4pt}] (3.4,3)--(5.6,3);\draw (5.6,3)--(6,3);
\draw (0,3)--(0,0)--(2.6,0);\draw (2.6,-.15) rectangle (3.4,.15);\draw (3.4,0)--(6,0)--(6,3);\node[below] at (3,-.2){$R$};
\draw (0,3)--(0,5)--(1.3,5);\draw[accent,thick] (1.3,4.65)--(1.3,5.35);\draw[accent,thick] (1.6,4.8)--(1.6,5.2);\draw (1.6,5)--(3,5)--(3,3);
\node[above] at (1.45,5.4){$U_0$};\node[above] at (1.5,3.25){$N_1=300$};\node[above] at (4.7,3.25){$N_2=1700$};
\foreach \x/\lab in {0/A,3/B,6/C}{\fill (\x,3) circle (2pt);\node[below] at (\x,2.85){$\lab$};}
\draw[arr,accent] (.8,2.2)--(2.2,2.2) node[midway,below]{$i_1$};\draw[arr,main] (5.2,2.2)--(3.8,2.2) node[midway,below]{$i_2$};
\node[note] at (3,-1.4){两段共享同一磁通；电阻跨全绕组两端};
\end{tikzpicture}

\begin{tikzpicture}
\node[box] (n0) at (0,0) {短段正安匝\\$N_1i_1$};
\node[box] (n1) at (5,0) {长段反安匝\\$-N_2i_2$};
\node[box] (n2) at (5,-2.3) {共同磁通\\$\Phi=\mathcal P(N_1i_1-N_2i_2)$};
\node[box] (n3) at (0,-2.3) {磁场能\\$E_B=L_1i_{\rm mag}^2/2$};
\draw[arr] (n0)--(n1);\draw[arr] (n1)--(n2);\draw[arr] (n2)--(n3);
\node[note] at (2.5,-4.1) {$i_1=i_{\rm mag}+(N_2/N_1)i_2$；负载反映项不改变净安匝数};
\end{tikzpicture}

\begin{tikzpicture}
\draw[arr] (0,0)--(6,0) node[right]{$t/\mathrm s$};
\draw[arr] (0,0)--(0,4) node[above]{$i/\mathrm A$};
\draw[accent,thick] (0.00000,0.32583) -- (0.10204,0.38842) -- (0.20408,0.45100) -- (0.30612,0.51359) -- (0.40816,0.57617) -- (0.51020,0.63876) -- (0.61224,0.70134) -- (0.71429,0.76393) -- (0.81633,0.82651) -- (0.91837,0.88910) -- (1.02041,0.95168) -- (1.12245,1.01427) -- (1.22449,1.07685) -- (1.32653,1.13944) -- (1.42857,1.20202) -- (1.53061,1.26461) -- (1.63265,1.32719) -- (1.73469,1.38978) -- (1.83673,1.45236) -- (1.93878,1.51495) -- (2.04082,1.57753) -- (2.14286,1.64012) -- (2.24490,1.70270) -- (2.34694,1.76529) -- (2.44898,1.82787) -- (2.55102,1.89046) -- (2.65306,1.95304) -- (2.75510,2.01563) -- (2.85714,2.07821) -- (2.95918,2.14080) -- (3.06122,2.20338) -- (3.16327,2.26597) -- (3.26531,2.32855) -- (3.36735,2.39114) -- (3.46939,2.45372) -- (3.57143,2.51631) -- (3.67347,2.57889) -- (3.77551,2.64148) -- (3.87755,2.70406) -- (3.97959,2.76665) -- (4.08163,2.82923) -- (4.18367,2.89182) -- (4.28571,2.95440) -- (4.38776,3.01699) -- (4.48980,3.07957) -- (4.59184,3.14216) -- (4.69388,3.20474) -- (4.79592,3.26733) -- (4.89796,3.32991) -- (5.00000,3.39250);
\draw[main,dashed] (0.00000,0.00000) -- (0.10204,0.06259) -- (0.20408,0.12517) -- (0.30612,0.18776) -- (0.40816,0.25034) -- (0.51020,0.31293) -- (0.61224,0.37551) -- (0.71429,0.43810) -- (0.81633,0.50068) -- (0.91837,0.56327) -- (1.02041,0.62585) -- (1.12245,0.68844) -- (1.22449,0.75102) -- (1.32653,0.81361) -- (1.42857,0.87619) -- (1.53061,0.93878) -- (1.63265,1.00136) -- (1.73469,1.06395) -- (1.83673,1.12653) -- (1.93878,1.18912) -- (2.04082,1.25170) -- (2.14286,1.31429) -- (2.24490,1.37687) -- (2.34694,1.43946) -- (2.44898,1.50204) -- (2.55102,1.56463) -- (2.65306,1.62721) -- (2.75510,1.68980) -- (2.85714,1.75238) -- (2.95918,1.81497) -- (3.06122,1.87755) -- (3.16327,1.94014) -- (3.26531,2.00272) -- (3.36735,2.06531) -- (3.46939,2.12789) -- (3.57143,2.19048) -- (3.67347,2.25306) -- (3.77551,2.31565) -- (3.87755,2.37823) -- (3.97959,2.44082) -- (4.08163,2.50340) -- (4.18367,2.56599) -- (4.28571,2.62857) -- (4.38776,2.69116) -- (4.48980,2.75374) -- (4.59184,2.81633) -- (4.69388,2.87891) -- (4.79592,2.94150) -- (4.89796,3.00408) -- (5.00000,3.06667);
\draw[main,thick] (0.00000,0.05750) -- (0.10204,0.05750) -- (0.20408,0.05750) -- (0.30612,0.05750) -- (0.40816,0.05750) -- (0.51020,0.05750) -- (0.61224,0.05750) -- (0.71429,0.05750) -- (0.81633,0.05750) -- (0.91837,0.05750) -- (1.02041,0.05750) -- (1.12245,0.05750) -- (1.22449,0.05750) -- (1.32653,0.05750) -- (1.42857,0.05750) -- (1.53061,0.05750) -- (1.63265,0.05750) -- (1.73469,0.05750) -- (1.83673,0.05750) -- (1.93878,0.05750) -- (2.04082,0.05750) -- (2.14286,0.05750) -- (2.24490,0.05750) -- (2.34694,0.05750) -- (2.44898,0.05750) -- (2.55102,0.05750) -- (2.65306,0.05750) -- (2.75510,0.05750) -- (2.85714,0.05750) -- (2.95918,0.05750) -- (3.06122,0.05750) -- (3.16327,0.05750) -- (3.26531,0.05750) -- (3.36735,0.05750) -- (3.46939,0.05750) -- (3.57143,0.05750) -- (3.67347,0.05750) -- (3.77551,0.05750) -- (3.87755,0.05750) -- (3.97959,0.05750) -- (4.08163,0.05750) -- (4.18367,0.05750) -- (4.28571,0.05750) -- (4.38776,0.05750) -- (4.48980,0.05750) -- (4.59184,0.05750) -- (4.69388,0.05750) -- (4.79592,0.05750) -- (4.89796,0.05750) -- (5.00000,0.05750);
\node[above] at (4.1,3.2){$i_1$};\node[below right] at (3.7,2.1){$i_{\rm mag}$};\node[right] at (5,.15){$i_2$};\node[below] at (2.5,0){$.10$};\node[below] at (5,0){$.20$};\node[left] at (0,.326){$.283$};\node[note] at (2.8,-1.2){理想接通后：长段恒定，短段含负载反映初值};
\end{tikzpicture}

\begin{tikzpicture}
\node[box] (n0) at (0,0) {输入功率\\$U_0(i_1+i_2)=2.50\,\mathrm W$};
\node[box] (n1) at (5,0) {磁场储能率\\$U_0i_{\rm mag}=2.00\,\mathrm W$};
\node[box] (n2) at (5,-2.3) {电阻发热\\$Ri_2^2=0.500\,\mathrm W$};
\node[box] (n3) at (0,-2.3) {能量闭合\\$P_b=\dot E_B+P_R$};
\draw[arr] (n0)--(n1);\draw[arr] (n1)--(n2);\draw[arr] (n2)--(n3);
\node[note] at (2.5,-4.1) {数值对应 $t=0.100\,\mathrm s$；电池电流是两个支路电流之和};
\end{tikzpicture}

\begin{tikzpicture}
\node[box] (n0) at (0,0) {短段电压\\$N_1\dot\Phi=U_0$};
\node[box] (n1) at (5,0) {全绕组电压\\$N\dot\Phi=Ri_2$};
\node[box] (n2) at (5,-2.3) {安匝平衡\\$i_1=i_{\rm mag}+N_2i_2/N_1$};
\node[box] (n3) at (0,-2.3) {题问电流\\$1.61667\,\mathrm A,\ 0.0500\,\mathrm A$};
\draw[arr] (n0)--(n1);\draw[arr] (n1)--(n2);\draw[arr] (n2)--(n3);
\node[note] at (2.5,-4.1) {完全耦合初值满足净安匝数连续，逐支路电流不必连续};
\end{tikzpicture}
\end{document}
